% !TeX root = ../main.tex
\clearpage
\phantomsection\label{epilogue}\pdfbookmark[0]{Epilogue}{epilogue}
\vspace*{5mm}
{\sffamily\bfseries\color{BookBlue}\fontsize{26.25}{34}\selectfont\noindent Epilogue\par}
\vspace{20pt}
\fontsize{12}{23.4}\selectfont
\setlength{\parindent}{2\ccwd}

写这本手册的原因，起于一次大额报销困难的经历以及被拖欠报销导致部分发票逾期不可报销，两件事的叠加使得我对于
体育部的报销流程产生了困惑，其低效率之低令人叹为观止。

我会这么认为是因为我同时负责我个人的大创项目的报销。大创的报销是由我所在的学院（人工智能学院）负责的。我们学院
授权学生可以直接作为制单人提交报销单据，全程线上提交审批，经学院内部负责人审批后到项目负责人最后移交财务，
有缺失材料或者补充信息也是直接在线上提交或修改。从制单提交以及中间的审批和材料补充到最后下账，不会超过一周。
这样对比下来比赛相关的报销就显得很慢了。

当然我也理解，外出比赛涉及的报销手续和需要的材料比一般学生大创需要的材料要复杂得多，涉及到的人员也更多，
整理材料的时间更长和难度也更大，但是这也不应该成为拖延得如此夸张，以至于出现了逾期坏账的情况。

\begin{tcolorbox}[handbook,colback=white,borderline west={1.5pt}{0pt}{BookPurple},left=11pt,top=1pt,bottom=1pt,before skip=13pt,after skip=13pt,fontupper=\sffamily\bfseries\fontsize{11.75}{22}]
大家自己垫付费用代表学校参加比赛，拿回属于自己的报销款项是理所应当的事情，而不是需要乞讨才能获得的施舍。
\end{tcolorbox}

我们无法去改变体育部的保留流程，部分材料和项目授权我们无能为力。我们能做的，便是将队员个人的才来整理得尽可能
完备整齐，构建一个清晰严格的报销材料提交体系，减少不必要的来回询问和材料补充，尽量让报销流程顺畅。

\vspace{6pt}\begin{flushright}\fontsize{10.5}{18}\selectfont
\Strong{PhSeCl}\\
2026/11/05\\
at SYSU Zhuhai Campus
\end{flushright}
